\documentclass[11pt]{article}
\usepackage[margin=0.78in]{geometry}
\usepackage{amsmath,amssymb,bm}
\usepackage{booktabs}
\usepackage{enumitem}
\usepackage{hyperref}
\hypersetup{colorlinks=true, linkcolor=black, urlcolor=black}
\setlength{\parindent}{0pt}
\setlength{\parskip}{0.55em}

\title{Cyrus Guided Notes: MIT 6.241J Dynamic Systems and Control}
\author{MIT course pack}
\date{2026-05-15}

\begin{document}
\maketitle

\section*{Course source}
\textbf{Official source:} \url{https://ocw.mit.edu/courses/6-241j-dynamic-systems-and-control-spring-2011/}

\textbf{Why this course is in the system.} Core control course for state space, controllability, observability, stability, and feedback.

\textbf{Learning output.} Priority course for a complete guided PDF connected to existing control formula visuals.

\section*{Prerequisite checkpoint}
\begin{itemize}[leftmargin=1.4em]
  \item Linear algebra
  \item Eigenvalues
  \item Differential equations
\end{itemize}

\section*{Formula checkpoint}
Do not memorize these first. Read each formula as a sentence: what changes, what is observed, what is optimized, and what output it produces.
\begin{align*}
  \dot{\bm{x}}=A\bm{x}+B\bm{u}\\
  e^{At}=\mathcal{L}^{-1}\{(sI-A)^{-1}\}\\
  \operatorname{rank}[B\;AB\;\cdots\;A^{n-1}B]=n\\
  \bm{u}=-K\bm{x}
\end{align*}

\section*{Visual page}
Make state propagation visible: A moves the state, B injects input, K closes feedback.

\section*{GoodNotes pages}
\begin{itemize}[leftmargin=1.4em]
  \item Page MIT-C001: state space
  \item Page MIT-C002: controllability
  \item Page MIT-C003: feedback
\end{itemize}

\section*{Web interaction}
A two-state system where dragging eigenvalues changes stability.

\section*{Obsidian and Notion}
\begin{tabular}{p{0.22\linewidth}p{0.68\linewidth}}
\toprule
Layer & Output \\
\midrule
Obsidian & MIT Control -> 6.241J dynamic systems \\
Notion & Control gate status and solved derivation evidence. \\
\bottomrule
\end{tabular}

\section*{First study move}
Open the official source, copy the prerequisite checkpoint into GoodNotes, then write one formula in your own words before watching or reading more.

\end{document}
